\documentclass[11pt]{article}

% ----------- Font Packages ------------
\usepackage[T1]{fontenc}
\usepackage[utf8]{inputenc}

% ----------- Math Formatting ----------
\usepackage{amsmath, amssymb, amsthm}
\usepackage{dsfont}
\usepackage[most]{tcolorbox}
\tcbuselibrary{theorems}

% Cross-referencing and hyperlinks
\usepackage{hyperref}
\usepackage{cleveref}

% ----------- Layout and Spacing -------
\usepackage[letterpaper, top=0.75in, bottom=0.75in, left=1in, right=1in, heightrounded]{geometry}
\usepackage{setspace}
\onehalfspacing

% ----------- Misc Enhancements --------
\usepackage{microtype}
\usepackage{enumitem}
\usepackage{fancyhdr}
\pagestyle{fancy}
\fancyhead[L]{Homework Three}

\usepackage{bookmark}

\hypersetup{
    colorlinks=true,
    linkcolor=blue,
    urlcolor=blue,
    citecolor=blue
}

\title{Algebra Homework: Laws of Exponents and Roots}
\author{Name: \underline{\hspace{3cm}}}
\date{Due: Nov 17}

\begin{document}

\maketitle

% ---------------------- Foundational ----------------------
\section*{Foundational Questions}

\begin{itemize}
    \item What is the relationship between squaring a number and taking its square root? Give an example using both operations on the same number.
    \item Explain why \(\sqrt{a^2} = |a|\) for any real number \(a\). Why can’t we simply say it equals \(a\)?
    \item What is the difference between a rational exponent such as \(x^{1/3}\) and a negative exponent such as \(x^{-3}\)?
\end{itemize}

% ---------------------- Computation ----------------------
\section*{Exponent and Radical Practice}

Simplify each expression. Show all steps, and express radicals using exponents where helpful.

\begin{enumerate}[label=\textbf{(\arabic*)}, itemsep=1em]
    \item \(\sqrt{49}\)
    \item \(\sqrt{75} \text{ (simplify completely)}\)
    \item \( \sqrt[3]{216} \)
    \item \( 16^{3/4} \)
    \item \( (9x^2)^{1/2} \)
    \item \( (25y^4)^{1/2} \)
    \item \( 81^{5/4} \)
    \item Convert to radical form: \(x^{7/2}\)
    \item Convert to exponent form: \(\sqrt[5]{a^3}\)
    \item Simplify: \(\dfrac{4x^{3/2}}{2x^{1/2}}\)
    \item Simplify: \((3x^{-2})(6x^5)\)
    \item \textbf{Challenge:} Simplify completely:
    \[
        \dfrac{(64x^6)^{1/3}}{(4x^{-1})^{1/2}}.
    \]
\end{enumerate}

% ---------------------- Word Problems ----------------------
\section*{Word Problems}

\begin{enumerate}[label=\textbf{(\arabic*)}, itemsep=1.5em]
    \item The area of a square is \(144\) square centimeters.  
    What is the length of one side?

    \item A cube has a volume of \(3.43 \times 10^{-3} \text{ m}^3\).  
    What is the side length of the cube? (Hint: use cube roots.)

    \item A bacteria culture grows according to the function  
    \[
       P(t) = 5.0 \times 10^4 \cdot (2^{t/3}),
    \]
    where \(t\) is in hours.  
    Find the population after 9 hours.

    \item The intensity of a sound (in watts per square meter) can be written as  
    \[
       I = 10^{-12} \cdot 10^{L/10},
    \]
    where \(L\) is the decibel level.  
    Compute the intensity when \(L = 70\).

    \item The gravitational force between two masses is given by  
    \[
        F = G \dfrac{m_1 m_2}{r^2}.
    \]
    If \(G = 6.67 \times 10^{-11}\), \(m_1 = 4.0 \times 10^3\),  
    \(m_2 = 2.5 \times 10^{4}\), and \(r = 5.0 \times 10^2\), compute the force in scientific notation.

    \item \textbf{Challenge Problem}  
    The brightness of a star is proportional to the inverse square of its distance from Earth:
    \[
       B = \dfrac{k}{d^2}.
    \]
    Suppose Star A is \(4 \times 10^{17}\) meters away and Star B is \(1 \times 10^{18}\) meters away.  
    \begin{enumerate}[label=\textbf{(\alph*)}]
        \item Write the ratio \(\dfrac{B_A}{B_B}\) in terms of their distances.
        \item Simplify the ratio using exponent rules.
        \item Which star appears brighter, and by how many times?
    \end{enumerate}
\end{enumerate}

\bigskip

\noindent\textbf{Extra Practice (optional):}
\begin{enumerate}[label=\textbf{(\alph*)}, itemsep=0.5em]
    \item Simplify: \((8x^{-3/2}) (2x^{5/2})\)
    \item Simplify: \(\dfrac{(3 \times 10^{-4})^{3/2}}{(3 \times 10^{-2})^{1/2}}\)
    \item Evaluate: \(\sqrt[4]{16^3}\)
\end{enumerate}

\end{document}
